\subsection{Healthcare \& Medicine Agents}
\label{sec:app-healthcare}



Healthcare and medical agents seek to support the full clinical decision pipeline, from initial symptom triage to treatment planning and integrating LLMs with structured patient records, medical ontologies and expert guidelines. Unlike general assistants, these systems must operate under strict safety constraints, multi-modal evidence and legal justification.

In this subsection, we begin with the foundational layer (Section~\ref{sec:app-healthcare-foundational}), which includes medical and diagnostic reasoning and tool-augmented access to various biomedical knowledge bases and APIs. Building on these primitives, the self-evolving layer (Section~\ref{sec:app-healthcare-evolve}) examines memory, feedback and reflective modules that allow these agents to accumulate patient-specific context, adapt to longitudinal trajectories and revise clinical plans over time. Finally, the collective layer (Section~\ref{sec:app-healthcare-multi}) highlights multi-agent collaboration, which includes doctor–agent co-planning, human–AI shared autonomy and specialist model ensembles.



\subsubsection{Foundational agentic reasoning}
\label{sec:app-healthcare-foundational}



\paragraph{Planning.}


Planning is a core capability for healthcare agents, which enables them to structure long-horizon clinical pathways into diagnostic and treatment phases, refine workflows dynamically as patient conditions evolve and coordinate across teams and tools toward cohesive care delivery. We discuss several various recent advancements as follows. For instance, a recent agentic clinical system~\cite{ferber2024autonomousartificialintelligenceagents} orchestrates specialized tools and guideline citations to support oncology decision-making, EHRAgent \cite{wenqi2024ehragent} decomposes multi-table EHR inference into code-execution steps with feedback learning and PathFinder \cite{ghezloo2025pathfinder} presents a multi-agent, multi-modal histopathology workflow for diagnostic reasoning.  

Other frameworks model planning as an explicit orchestration layer across levels of abstraction. For example, MedAgent‑Pro \cite{wang2025medagent} proposes a hierarchical workflow which first generates disease-level diagnostic plans from guideline criteria and then dispatches tool-agent modules for execution. MedOrch \cite{he2025medorchmedicaldiagnosistoolaugmented} treats tool invocation itself as a planning primitive across modalities, orchestrating reasoning agents for multi-step diagnostic execution. On the other hand, ClinicalAgent \cite{yue2024clinicalagentclinicaltrialmultiagent} coordinates multi-agent workflows for clinical planning, leveraging LLM reasoning to allocate tools and synthesize evidence. In addition, planning in healthcare agents is increasingly adaptive, responding to new information and evolving contexts. For example, DoctorAgent‑RL \cite{feng2025doctoragentrlmultiagentcollaborativereinforcement} models clinical consultation as a dynamic decision-making process under uncertainty, optimizing questioning strategies and diagnostic paths via reinforcement learning; while DynamiCare \cite{shang2025dynamicaredynamicmultiagentframework} adjusts specialist-agent teams across multi-round interactions as new patient information emerges.



\paragraph{Tool-use.}


Tool integration significantly expands a healthcare agent’s action space, enabling precise calculations, medical image interpretation and access to specialized databases. Recent studies are summarized as follows. Several systems explicitly foreground extensibility. MedOrch~\cite{he2025medorchmedicaldiagnosistoolaugmented} introduces a modular architecture that allows new diagnostic APIs to be incorporated without retraining, while TxAgent~\cite{gao2025txagentaiagenttherapeutic} integrates over two hundreds pharmacological tools to support therapeutic decision-making across drug–disease–treatment relationships. AgentMD~\cite{jin2024agentmdempoweringlanguageagents} similarly curates and leverages over two thousands executable clinical calculators to learn risk-prediction pipelines.

Other approaches focus on structured function calling for safe execution. For example, LLM-based agents can reliably invoke bedside calculators when provided with explicit function signatures, ensuring arithmetic correctness in dosing and risk scoring~\cite{Goodell_2025}. MeNTi~\cite{zhu2025mentibridgingmedicalcalculator} goes further by enabling nested tool calls across multi-step medical calculators. Complementing these text-based integrations, MMedAgent~\cite{li2024mmedagent} demonstrates that agents can learn to select among multi-modal tools. 

In addition, execution is crucial for translating high-level clinical plans into concrete actions such as code operations, database queries or robotic procedures. VoxelPrompt~\cite{hoopes2024voxelpromptvisionlanguageagentgrounded} embed 3-D volumetric priors so that language instructions drive spatial segmentation and analysis of medical image volumes. On the other hand, embodied ultrasound-robot controllers~\cite{xu2024enhancingsurgicalrobotsembodied} translate LLM-generated plans into closed-loop robotic scanning via a “think-observe-execute” loop.
Adaptive reasoning-and-acting systems~\cite{dutta2024adaptivereasoningactingmedical} further refine both the reasoning and actions over time in simulated clinical environments. In medical imaging, systems like MedRAX~\cite{fallahpour2025medraxmedicalreasoningagent} materializes multi-step reasoning by integrating specialized chest-X-ray tools and LLM reasoning into an end-to-end diagnostic agent. PathFinder~\cite{ghezloo2025pathfinder} similarly executes multi-agent, multi-modal diagnostic workflows in histopathology.

Another class of healthcare agents deploys code-level workflows. For example, Conversational Health Agents~\cite{abbasian2024conversationalhealthagentspersonalized} compile dialogue actions into function calls and code execution for downstream processing, while EHRAgent~\cite{wenqi2024ehragent} materializes EHR operations via executable code. MedAgentGym~\cite{xu2025medagentgymscalableagentictraining} trains agents to produce code that is directly executed and graded, enforcing reliability of reasoning traces. DoctorAgent-RL~\cite{feng2025doctoragentrlmultiagentcollaborativereinforcement} validates multi-turn dialogue acts by executing reinforcement-learned strategies in simulated consultations, while AIPatient~\cite{yu2025simulated} materializes realistic patient scenarios for execution-based evaluation and another recent study~\cite{almansoori2025selfevolvingmultiagentsimulationsrealistic} demonstrates how self-evolving multi-agent simulations allow execution behaviors themselves to improve over time.



\paragraph{Search and retrieval.}
Search-based agents enhance clinical decision-making by linking LLM reasoning with external biomedical knowledge sources. For instance, MeNTi~\cite{zhu2025mentibridgingmedicalcalculator} supplements therapeutic reasoning by bridging LLM calls into multi-step medical calculators while EHRAgent~\cite{wenqi2024ehragent} dynamically executes code operations over multi-table EHR data to support complex tabular inference. Conversational Health Agents~\cite{abbasian2024conversationalhealthagentspersonalized} enrich personalized dialogue by integrating developer-defined external sources and orchestrating action flows. Another line of work explicitly embeds retrieval-augmented generation (RAG) into healthcare agents. For example, CLADD~\cite{lee2025ragenhancedcollaborativellmagents} retrieves molecular graphs and prior assay results before proposing compound hypotheses and MedReason~\cite{wu2025medreasonelicitingfactualmedical} issues targeted knowledge-graph sub-queries to anchor each reasoning step for clinical QA.




\subsubsection{Self-evolving agentic reasoning}
\label{sec:app-healthcare-evolve}



Self-evolving capabilities enable healthcare agents to maintain longitudinal clinical coherence. Representative use cases include accumulating relevant medical context across encounters, updating beliefs as new evidence arrives and revising decisions when inconsistencies surface. In the following paragraphs, we examine how memory, feedback and reflective mechanisms collectively turn clinical reasoning from a one-shot prediction into a continually improving care process.


\paragraph{Memory.}
Persistent memory is essential for tracking medical or patient history and maintaining context across interactions. For instance, epidemic-modeling agents~\cite{williams2023epidemicmodelinggenerativeagents} maintain temporal contact histories to trace infection chains over time; while MedAgentSim~\cite{almansoori2025selfevolvingmultiagentsimulationsrealistic} stores experience histories and refine diagnostic strategies over time. In structured data settings, EHRAgent~\cite{wenqi2024ehragent} records intermediate computations over tabular EHRs so subsequent steps can reference prior results. EvoPatient~\cite{du2025llmssimulatestandardizedpatients} interleaves memory with coevolution maintains evolving clinical state across dialogue phases while AIPatient~\cite{yu2025simulated} persists longitudinal EHR-derived variables to drive consistent responses. Multi-agent systems such as MedOrch~\cite{he2025medorchmedicaldiagnosistoolaugmented} contain clinical knowledge graph agent which can be considered as external memory that can be queried to retrieve known relationships or diagnostic patterns. 



\paragraph{Agentic Feedback and Reflection.}
Agentic feedback and self-reflection are complementary mechanisms that improve reliability and adaptability of healthcare agents. Feedback converts execution outcomes into learning signals: DynamiCare~\cite{shang2025dynamicaredynamicmultiagentframework} updates multi-agent treatment strategies when newly observed patient state contradicts prior plans; DoctorAgent-RL~\cite{feng2025doctoragentrlmultiagentcollaborativereinforcement} optimizes questioning policies from consultation rewards; and MedAgentGym~\cite{xu2025medagentgymscalableagentictraining} enforces correctness by executing and grading generated code. Tool-use pipelines also propagate execution feedback. For example, the success/failure of table queries in EHRAgent~\cite{wenqi2024ehragent} or calculator calls in MeNTi~\cite{zhu2025mentibridgingmedicalcalculator} and clinical-calculation agents~\cite{Goodell_2025} to refine subsequent actions.





\subsubsection{Collective multi-agent reasoning}
\label{sec:app-healthcare-multi}


Multi-agent collaboration is central to healthcare AI, since clinical decision-making often depends on consensus among specialists, negotiation of competing hypotheses and coordination across roles such as physicians, patients and trial designers. In the following, we discuss several strands of research centered around multi-agent capabilities.



For collaborative decision-making, notable frameworks include MDAgents~\cite{kim2024mdagents}, which automatically assigns tailored collaboration structures to teams of LLMs depending on medical task complexity, and DoctorAgent-RL~\cite{feng2025doctoragentrlmultiagentcollaborativereinforcement}, which uses a multi-agent reinforcement-learning framework to optimize multi-turn doctor-patient consultation dialogues. In addition, Agent-derived Multi-Specialist Consultation (AMSC)~\cite{wang2024directdiagnosisllmbasedmultispecialist} explores staged multi-specialist dialogues for differential diagnosis that mimics the medical scene of a patient consulting with multiple specialists. Other notable works include ClinicalAgent~\cite{yue2024clinicalagentclinicaltrialmultiagent}, which organizes clinical trial workflows via role-based agent collaboration / LLM reasoning and PathFinder~\cite{ghezloo2025pathfinder}, which integrates a diverse set of agents that can gather evidence and provide comprehensive diagnoses with natural language explanations. 

On the other hand, there are studies focusing on simulation-driven collaboration. These works highlight how multi-agent setups enrich training and evaluation. MedAgentSim~\cite{almansoori2025selfevolvingmultiagentsimulationsrealistic} co-evolves doctor and patient agents to simulate real-world multi-turn clinical interactions, and EvoPatient~\cite{du2025llmssimulatestandardizedpatients} uses co-evolution of patient and doctor agents to generate diagnostic dialogue data and therefore gathers experience to improve the quality of both questions and answers to enable accurate human doctor training. In addition, DynamiCare~\cite{shang2025dynamicaredynamicmultiagentframework} initiates a team of specialist agents that iteratively queries the patient system to integrate new information and adapts the composition and strategy. Finally, medical agents can also collaborative to aid medical reasoning process. For example, MedAgents~\cite{xiangru2023medagents} demonstrates zero-shot cooperation among domain-specialist agents in medical reasoning tasks, CLADD~\cite{lee2025ragenhancedcollaborativellmagents} uses retrieval-augmented generation to support drug-discovery workflows across agents and GMAI-VL-R1~\cite{su2025gmaivlr1harnessingreinforcementlearning} combines multi-modal reasoning and reinforcement learning in a multi-agent framework to support large-scale medical decision-making.


